%%%%%%%%%%%%%%%%%%%%%%%%%%%%%%%%%%%%%%%%%%%%%%%%%%%%%%%%%%%%%%%%%%%%%%%%%%%
%
% Generic template for TFC/TFM/TFG/Tesis
%
% $Id: resumen.tex,v 1.9 2015/06/05 00:10:31 macias Exp $
%
% By:
%  + Javier Macías-Guarasa. 
%    Departamento de Electrónica
%    Universidad de Alcalá
%  + Roberto Barra-Chicote. 
%    Departamento de Ingeniería Electrónica
%    Universidad Politécnica de Madrid   
% 
% Based on original sources by Roberto Barra, Manuel Ocaña, Jesús Nuevo,
% Pedro Revenga, Fernando Herránz and Noelia Hernández. Thanks a lot to
% all of them, and to the many anonymous contributors found (thanks to
% google) that provided help in setting all this up.
%
% See also the additionalContributors.txt file to check the name of
% additional contributors to this work.
%
% If you think you can add pieces of relevant/useful examples,
% improvements, please contact us at (macias@depeca.uah.es)
%
% You can freely use this template and please contribute with
% comments or suggestions!!!
%
%%%%%%%%%%%%%%%%%%%%%%%%%%%%%%%%%%%%%%%%%%%%%%%%%%%%%%%%%%%%%%%%%%%%%%%%%%%

\chapter*{Resumen}
\label{cha:resumen}
\markboth{Resumen}{Resumen}

\addcontentsline{toc}{chapter}{Resumen}

El Problema del Viajante es un problema de optimización combinatoria NP-completo de gran relevancia teórica y
práctica, para el que las metaheurísticas ofrecen soluciones de alta calidad en tiempos de ejecución razonables. Este
trabajo implementa, desde cero y en Python, cuatro metaheurísticas clásicas: algoritmos genéticos, colonia de hormigas,
recocido simulado y búsqueda tabú, sobre una representación común del problema, y evalúa su rendimiento mediante un 
experimento que replica el diseño de un estudio de referencia sobre un conjunto de instancias estandarizadas. Los 
hiperparámetros de cada algoritmo se seleccionan y justifican a partir de la literatura existente, en lugar de ajustarse de
forma específica para cada una las instancias evaluadas. Se analiza el comportamiento comparativo de las cuatro técnicas
atendiendo a la calidad de las soluciones obtenidas, el tiempo de ejecución y su robustez frente al tamaño del problema.

\textbf{Palabras clave:} \myThesisKeywords.

%%% Local Variables:
%%% TeX-master: "../book"
%%% End:


